\section{The actual original affine charts of the native rank-one valued point}

\begin{theorem}[All original chart-ring values agree on the actual native point]
\label{thm:fc-native-rank-one-analytic-charts}
\leanok
\lean{EnvelopeProperty.MultiModel.nativeRankOneHomogeneousSectionsIso, EnvelopeProperty.MultiModel.nativeRankOneOriginalSectionsIso, EnvelopeProperty.MultiModel.nativeRankOneAbsoluteChartHomeomorph, EnvelopeProperty.MultiModel.nativeRankOneAbsoluteChartHomeomorph_val_apply, EnvelopeProperty.MultiModel.nativeRankOne_chartMono_eq_pow, EnvelopeProperty.MultiModel.nativeRankOnePoint_eval_originalChart}
\uses{thm:fc-native-rank-one-analytic-point, thm:fc-native-rank-one-chart-comparison, lem:gl-charts}
The genuine section ring of each original nonzero native coordinate
open is isomorphic to the SAME original chart algebra. Its entire
absolute native valued-point chart is the actual original affine
seminorm spectrum. For every original analytic point in that chart,
evaluation of EVERY corresponding true chart section is exactly the
SAME original affine chart seminorm value.
\end{theorem}
\begin{proof}
Compose the inverse of the proved original chart-ring isomorphism with
the actual native homogeneous fraction-to-section isomorphism. The
genuine affine-open valued-point theorem and the resulting whole-ring
isomorphism give the whole absolute chart homeomorphism, preserving
each section value. Every original chart element is an original
same-degree coefficient/coordinate fraction. Its matching actual
homogeneous fraction maps to it under the whole proved chart-ring
isomorphism. The true native valued point evaluates this fraction by
the original numerator value divided by the coordinate-power value.
In rank one the original chart denominator monomial is that SAME
coordinate power. The original chart seminorm has exactly this same
fraction formula, proving the equality for every original chart
element.
\end{proof}
